\documentclass[11pt,a4paper]{ctexart}
\usepackage[margin=23mm]{geometry}
\usepackage{amsmath,amssymb,mathtools,enumitem,needspace}
\usepackage[hidelinks]{hyperref}
\setlength{\parskip}{0.35em}
\setlength{\emergencystretch}{2em}
\setlist[enumerate]{leftmargin=2em,itemsep=0.2em}
\allowdisplaybreaks[2]
\newcommand{\R}{\mathbb R}
\newcommand{\N}{\mathbb N}
\newcommand{\dd}{\,\mathrm d}
\newcommand{\supp}{\operatorname{supp}}
\newcommand{\source}[1]{\par\noindent\textbf{【#1】}\quad}
\newcommand{\result}[3]{\Needspace{6\baselineskip}\subsection*{#1\quad #2}\label{st:#1}\source{#3}}
\newcommand{\ex}[3]{\Needspace{7\baselineskip}\subsection*{#1\quad #2}\label{ex:#1}\noindent{\small\textbf{出处与处理：}#3}\par}
\newcommand{\scope}{本次范围为课程教材《高等微积分学》第4.1--4.4节（原书75--98页，PDF76--99页，仅取第98页上半的第4.4节结尾），不含第4.5节广义函数简介、试验函数空间及基本解。第3.5节多项式逼近、第5.2节Fej\'er定理只作为卷积逼近的关联结论。课堂lecture007中的积分渐近另列为课堂补充。}
